\begin{table}[!htbp]
\caption{Operational L/K/S/T definitions and illustrative Chinese spans under handbook v4.2.14.}
\label{tab:guide-types}
\centering\small
\setlength{\tabcolsep}{4pt}
\renewcommand{\arraystretch}{1.12}
\begin{tabularx}{\linewidth}{@{}P{0.19\linewidth}P{0.47\linewidth}L@{}}
\toprule
Label & Operational definition & Chinese examples and English glosses \\
\midrule
L: Language & Natural-language names, proficiency levels, examinations, and certificates. & \zh{英语六级}\par\emph{College English Test Band 6}\par\smallskip\zh{日语N2}\par\emph{Japanese proficiency at N2 level} \\[3pt]
K: Knowledge-related requirements & Facts, principles, theories, and fields of study; educational and non-language qualifications; and explicit industry-background requirements. & \zh{SQL原理}\par\emph{Principles of SQL}\par\smallskip\zh{本科及以上学历}\par\emph{Bachelor's degree or higher} \\[3pt]
S: Occupational skill & Occupational activities and the application of tools or methods. & \zh{使用}[\zh{SQL}]\par\emph{Use [SQL]}\par\smallskip\zh{数据分析}\par\emph{Data analysis} \\[3pt]
T: Transversal competency & General communication, cooperation, cognition, self-management, and dispositions. & \zh{客户汇报}\par\emph{Reporting to clients}\par\smallskip\zh{责任心}\par\emph{Sense of responsibility} \\
\bottomrule
\end{tabularx}
\end{table}
